\section{Revised Harness Validation}
\para{Setup.} On October 29 we tested the revised harness on the UCSD-hosted
\texttt{api-glm-5.3-flash} route through TritonAPI, with strict JSON output.
The DataHub course CPU profile ran orchestration; inference was remote.
The route is not an immutable checkpoint. Complete traces, responses, configurations,
Python/package versions and source snapshots accompany the results.

\para{Live pilot.} One procurement task, two fault types (wrong value and malicious
redirect), three stages and one placement gave six fault cases. Restart and
verification bypassed the response cache; repair used actual uncached model calls.
Median restart and verification times were 7.53\,s and 17.32\,s. These include network
and local instrumentation; only these baseline timings are live measurements.
The clean run passed, but only one of two uncached clean repeats passed.
Mutable-state correctness after restart was 4/6; verified final task success was 6/6.
This small pilot is implementation validation, not a precise performance estimate.

\para{Propagation validation.} Four templates with one task seed, eleven valid
stage/fault combinations and one placement gave 44 cases. Clean baseline success
was 3/4; eight uncached clean repeats passed 7/8. Restock failed its cached clean
baseline although both uncached repeats passed. We retain all 44 cases and
separately evaluate the 33 cases whose baseline completed cleanly:
\begin{center}\small
\begin{tabular}{@{}lc@{}}
\toprule
Policy & Final success\\
\midrule
No recovery & 11/33\\
Full restart & 29/33\\
Verify + actual retry & 31/33\\
\bottomrule
\end{tabular}
\end{center}
Always-verify detected the injected faults in all 44 cases but achieved final
task success in only 40/44, with no attacker-directed actions remaining.
Detection therefore does not guarantee correct repair. Restock inability is not
counted as attack resistance. Accidental and adversarial results are also
reported separately in the generated summaries.

\para{Controls and failures.} This second validation used cached replay and a
1.25\,s minimum spacing for live requests. Its timings are not comparable to the
live pilot as serving-performance measurements. A preceding four-worker batch
failed on four requests before completing any fault case; its manifest and
partial traces are retained, and it is excluded from success-rate denominators.
Changed prompts still resample; two clean repeats do not establish reliable
variance estimates. No result here measures implemented selective recovery.

\para{Resource accounting.} There were 1,093 real requests across probes, pilot,
the failed batch and serial validation. At the campus-listed rates of \$0.15/\$0.50
per million input/output tokens, completed usage cost is estimated at \$0.1761.
Four failed requests retain \$0.0141 in conservative reservations, giving
about \$0.1902 charged-or-reserved against a \$5 cap. These are token-rate estimates,
not an authoritative billing statement. Credentials stayed in DataHub environment
variables and are excluded from all exported artifacts.
